% 文档类型设置：A4纸，10磅字体，双栏（期刊常用格式）
\documentclass[a4paper,10pt,twocolumn]{article}
\usepackage[UTF8]{ctex}  % 中文支持
\usepackage{amsmath}     % 数学公式扩展
\usepackage{graphicx}    % 图形插入支持
\usepackage{booktabs}    % 专业表格格式
\usepackage{geometry}    % 页面边距设置
\usepackage{cite}        % 参考文献引用
\usepackage{float}       % 图片表格位置控制

% 页面设置：左右边距2.5cm，上下边距2.5cm
\geometry{left=2.5cm,right=2.5cm,top=2.5cm,bottom=2.5cm}

% 论文标题信息
\title{基于改进LQR的无人车路径跟踪控制研究}
\author{张三\quad 李四\\
        某某大学控制科学与工程学院}
\date{\today}

\begin{document}

\maketitle  % 生成标题

% 摘要部分
\begin{abstract}
本文提出了一种改进的线性二次调节器(LQR)控制方法...（摘要内容）...
\textbf{关键词：}路径跟踪，LQR控制，状态反馈，无人驾驶
\end{abstract}

% 第一节：引言
\section{引言}
随着无人驾驶技术的发展...（引言内容）...

% 第二节：问题描述
\section{问题描述}
\subsection{系统模型}
考虑如下车辆动力学模型：
\begin{align}  % 对齐方程环境
\dot{x} &= A x + B u \label{eq:state} \\  % 添加标签便于引用
y &= C x
\end{align}
其中状态向量$x \in \mathbb{R}^n$，控制输入$u \in \mathbb{R}^m$...

% 第三节：控制器设计
\section{控制器设计}
\subsection{改进LQR算法}
传统LQR代价函数：
\begin{equation}
J = \int_0^\infty (x^T Q x + u^T R u) dt \label{eq:cost}
\end{equation}
本文改进后的代价函数加入误差积分项：
\begin{equation}
J' = J + \int_0^\infty (e^T W e) dt
\end{equation}
其中$e = y_{ref} - y$为跟踪误差...

% 第四节：仿真结果
\section{仿真结果}
% 表格示例
\begin{table}[H]  % [H]表示严格位置控制
\caption{控制器性能对比} \label{tab:compare}
\centering
\begin{tabular}{ccc} \toprule
控制器类型 & 超调量(\%) & 稳定时间(s) \\ \midrule
传统LQR & 12.5 & 2.8 \\
本文方法 & \textbf{4.2} & \textbf{1.6} \\ \bottomrule
\end{tabular}
\end{table}

% 图片示例
\begin{figure}[H]
\centering
\includegraphics[width=0.9\linewidth]{response_curve.pdf}  % 图片宽度为栏宽的90%
\caption{系统响应曲线对比} \label{fig:response}
\end{figure}

% 第五节：结论
\section{结论}
仿真结果表明...（结论内容）...

% 参考文献（建议使用bibtex，此处展示手动录入方式）
\begin{thebibliography}{9}
\bibitem{lqr1} Anderson B D O, Moore J B. Optimal control: linear quadratic methods[M]. Courier Corporation, 2007.
\bibitem{auto2} 王田苗, 陶永. 我国工业机器人技术现状与产业化发展战略[J]. 机械工程学报, 2014, 50(9):1-13.
\end{thebibliography}

% 附录（如果需要）
\appendix
\section{稳定性证明}
根据Lyapunov稳定性理论...

\end{document}